\documentclass[12pt,a4paper,titlepage,final]{article}
\usepackage[czech]{babel}
\usepackage[utf8]{inputenc}
\usepackage[bookmarksopen,colorlinks,plainpages=false,urlcolor=blue,unicode]{hyperref}
\usepackage{url}
\usepackage{float}
\usepackage[dvipdf]{graphicx}
\usepackage[top=3.5cm, left=2.5cm, text={17cm, 24cm}, ignorefoot]{geometry}

\begin{document}
\newpage
\pagestyle{empty}

%%%%%%%%%%%%%%%%%%%%%%%%%%%%%%%%%%%%%%%%%%%%%%%%%%%%%%%%%%%%%%%%%%%%%%%%%%%%%%
% titulní strana
\def\author{Lukáš Vokráčko}
\def\email{xvokra00@stud.fit.vutbr.cz}
\def\projname{Semestrální projekt - řešení obvodů}
\begin{titlepage}

% \vspace*{1cm}
\begin{figure}[!h]
\centering
% \includegraphics[width=3.5cm,keepaspectratio]{fit-logo.eps}
\includegraphics[width=7cm,keepaspectratio]{logo.eps}

\end{figure}

\vfill

\begin{center}
\begin{Huge}
Pokročilé asemblery\\
\end{Huge}
\begin{Large}
Benchmark - Rozostření
\end{Large}
\end{center}

\vfill

\begin{center}
\begin{Large}
\today
\end{Large}
\end{center}

\vfill

\begin{flushleft}
\begin{large}
\begin{tabular}{ll}
Autor: & Lukáš Vokráčko, \url{xvokra00@stud.fit.vutbr.cz} \\
 & Fakulta Informačních Technologií \\
 & Vysoké Učení Technické v~Brně \\
\end{tabular}
\end{large}
\end{flushleft}
\end{titlepage}


%%%%%%%%%%%%%%%%%%%%%%%%%%%%%%%%%%%%%%%%%%%%%%%%%%%%%%%%%%%%%%%%%%%%%%%%%%%%%%
% obsah

\tableofcontents

%%%%%%%%%%%%%%%%%%%%%%%%%%%%%%%%%%%%%%%%%%%%%%%%%%%%%%%%%%%%%%%%%%%%%%%%%%%%%%
% textova zprava
\newpage
\pagestyle{plain}
\pagenumbering{arabic}
\setcounter{page}{1}
%%%%%%%%%%%%%%%%%%%%%%%%%%%%%%%%%%%%%%%%%%%%%%%%%%%%%%%%%%%%%%%%%%%%%%%%%%%%%%

\section{Zadání}
	Optimalizujte implementaci Gaussovského rozostření obrázku použitím SSE a MMX instrukcí.

\section{Rozbor problému}
	Ke správnému řešení rychlostních problémů patří analýza již existujícího kódu a vhodné zvolení metod k optimalizaci.
	A to takových, které přinesou největší užitek v poměru k vynaloženému úsilí.

	Při řešení tohoto konkrétního problému jsem zvolil způsoby popsané v následující kapitole.

\section{Optimalizace}
	\subsection{Implementované optimalizace}
		\subsubsection{Minimalizace podmíněných skoků}
			Počet skoků lze minimalizovat analýzou podmínek a jejich vhodným rozmístěním v kódu.
			Je výhodné umisťovat kód větve, která se provádí častěji, přímo do hlavního bloku kódu
			a provádět instrukci skoku pouze v případě nutnosti vykonat kód statisticky méně prováděné větve.
			\\\\
			Statistiky vyhodnocení podmínek pro vzorový obrázek:

			\paragraph{Vertikální rozostření}
			\begin{center}
			    \begin{tabular}{ | l | l | l | }
			    \hline
			    Podmínka & Vyhodnocení jako \texttt{true} & Relativní četnost vyhodnocení \\ \hline
			    y \textgreater= y\_start & 1384020 & 0.984160 \\ \hline
			    y + radius \textgreater= y\_start & 10395 & 0.007391 \\ \hline
			    1. p \textgreater 0 & 0 & 0.000000 \\ \hline
			    2. p \textgreater 0 & 0 & 0.000000 \\ \hline
			    \end{tabular}
			\end{center}

			\paragraph{Horizontální rozostření}
			\begin{center}
			    \begin{tabular}{ | l | l | l | }
			    \hline
			    Podmínka & Vyhodnocení jako \texttt{true} & Relativní četnost vyhodnocení \\ \hline
			    x \textgreater= x\_start & 1383938 & 0.989998 \\ \hline
			    x + radius \textgreater= x\_start & 6524 & 0.004666 \\ \hline
			    1. p \textgreater  0 & 0 & 0.000000 \\ \hline
			    2. p \textgreater  0 & 0 & 0.000000 \\ \hline
			    \end{tabular}
			\end{center}


		\subsubsection{Použití registrů pro ukládání hodnot}
			Po analýze algoritmu jsem vhodně zvolil proměnné, které jsou nejčastěji používané
			a pro jejich uložení jsem použil vhodné registry. Tato čísla jsem získal použitím referenčního
			řešení na vzorovém obrázku.

			\begin{center}
			    \begin{tabular}{ | l | l | l | l |}
			    \hline
			    Proměnná & Počet použití & Zvolený registr & Část výpočtu \\ \hline
			    y & 7217556 & edx & vertical \\ \hline
			    x & 7217556 & edx & horizontal \\ \hline
			    radius & 5051598 & ecx & vertical, horizontal \\ \hline
			    dif & 2179031 & xmm7 & vertical, horizontal \\ \hline
			    sum & 2162882 & xmm6 & vertical, horizontal \\ \hline
			    p & 5778966 & xmm5 & vertical, horizontal \\ \hline
			    s & 1433206 & xmm4 & vertical, horizontal \\ \hline
			    w & 1433206 & xmm3 & vertical, horizontal \\ \hline
			    \end{tabular}
			\end{center}

		\subsubsection{Ukládání neměnných hodnot} \label{static}
			Některé hodnoty se v průběhu výpočtu nemění nebo se v každé iteraci mění pouze minimálně.
			Tyto hodnoty si lze vypočítat ještě před zahájením výpočtu a poté pracovat přímo s výsledky.

	\subsection{Další možné optimalizace}
		Dalšími možnými způsoby jak urychlit výpočet hodnot pro nový obrázek je několik.
		Vzhledem k nedostatku času jsem se je nerozhodl implementovat, pouze je zde zmíním.

		\subsubsection{Ukládání neměnných hodnot}
			Druh optimalizace zmíněný v \ref{static} je možné použít pro více hodnot,
			než jsem se rozhodl implementoval.
			Mezi tyto hodnoty patří následující:

			\begin{itemize}
				\item Vertikální rozostření
					\begin{itemize}
						\item \begin{verbatim}image[x + (y_start - radius - 1)*image_w]\end{verbatim}
						\item \begin{verbatim}image[x + (y_start - radius - 1)*image_w]\end{verbatim}
						\item \begin{verbatim}buffer[y+radius]\end{verbatim}
						\item \begin{verbatim}buffer[y-radius]\end{verbatim}
						\item \begin{verbatim}buffer[y]\end{verbatim}
					\end{itemize}
				\item Horizontální rozostření
					\begin{itemize}
						\item \begin{verbatim}new_image[x_start - radius - 1 + y*image_w]\end{verbatim}
						\item \begin{verbatim}buffer[x+radius]\end{verbatim}
						\item \begin{verbatim}buffer[x-radius]\end{verbatim}
						\item \begin{verbatim}buffer[x]\end{verbatim}
					\end{itemize}
			\end{itemize}

		\subsubsection{Paralelní zpracování}
			Další potenciální zrychlení je možné změnou indexování obrázků nebo cyklů tak,
			aby se hodnoty čtené/zapisované ve vnitřních cyklech vyskytovaly ve fyzické paměti za sebou.
			To by umožnilo plně využít potenciál MMX/SSE instrukcí a pracovat až s 16 hodnotami najednou.

			Tento druh optimalizace lze použít i pro \texttt{buffer} bez významnějších úprav
			zdrojového kódu algoritmu, nicméně tuto variantu se mi nepodařilo úspěšně\footnote{Se zachováním chyby v rozmezí 0-5}  implementovat.


\section{Závěr}
	Pomocí výše zmíněných optimalizací se mi podařilo výpočet urychlit $1.5\times$, nicméně se
	moje implementace neobešla bez chyb. Průměrná hodnota chyby na pixel činí $2.76$


\end{document}
